% ID: SEL-D0073-OB-E03
% ESTADO: cotejada. Segunda pasada visual y borrador compilado revisados.
% Fuente: Selectividad/sel_2008_matematicas_aplicadas/Mat CCSS 1 - Examen y criterios de correción.pdf
% SHA-256: 2461e302e686c60322df8edd0a004b11753e95c4824dc5835b16b76c05174e81
% Página: 2; opción B; ejercicio original 3.

\noindent\underline{Parte I}

De los 150 coches de un concesionario, 90 tienen motor diesel y el resto de gasolina. De los coches con motor diesel, 72 son nuevos y el resto usados; mientras que de los coches con motor de gasolina hay el mismo número de coches nuevos que de usados. Se elige, al azar, un coche de dicho concesionario; calcule la probabilidad de que:
\begin{enumerate}[label=\alph*)]
\item \textbf{(1 punto)} Sea nuevo.
\item \textbf{(1 punto)} Tenga motor diesel, sabiendo que es usado.
\end{enumerate}
\par\noindent\begin{minipage}{\linewidth}
\noindent\underline{Parte II}

\textbf{(2 puntos)} Una variable aleatoria sigue una ley Normal con desviación típica 6. \textquestiondown{}De qué tamaño, como mínimo, se debe elegir una muestra que nos permita estimar la media de esa variable con un error máximo de 2 y una confianza del $99\%$?
\end{minipage}\par
